\scanpage{6}
et $a$ peut être pris arbitrairement. On voit déjà
que le champ envisagé \uline{couvre}, et il reste
à \deleted{donner}\inserted{voir} la condition moyennant laquelle deux
objets sont tjrs loc isom, ou encore
pour deux éléments $a,b\in G(S)$ $(S\in\underline S)$,
$\varphi_a$ et $\varphi_b$ définissent des objets loc. isom.
Pour ceci, notons que les hom de
$(G,\varphi_a)$ dans $(G,\varphi_b)$ s'identifient aux $u\in G(S)$,
tels que
$(G,\varphi_a)\mathbin{\text{---}}$
\diagram{\begin{tikzcd}
G \arrow[r,no head,"\varphi_a"] \arrow[d,"\rho_u"'] & G \arrow[d,"\rho_u"]\\
G \arrow[r,no head,"\varphi_b"'] & G
\end{tikzcd}}
commute
où $\rho_u(g)=ug$, donc tels que
\[
u a f(g)=b f(ug)\qquad \forall g\in G(S'),\quad S'\text{ sur }S,
\]
ce qui s'écrit aussi simplement
\[
b^{-1}a=f(u)u^{-1}
\]
Donc on trouve que $\Sigma$ \uline{est une gerbe} sss
\uline{l'hom de faisceaux d'ensembles}
\[
u\longmapsto f(u)u^{-1}:G\longrightarrow G
\]
\uline{est un épimorphisme}. Alors $\unsure{\forall}$ objet de
$\Sigma$ est loc isom à $(G,\varphi_e=f)$ lui-$=$,
dont le faisceau des autom. est formé du
sous faisceau des $u$ de $G$ tels que $f(u)u^{-1}=\unsure{11^{-1}|ee^{-1}}=1$,
\mspage{5}
